\documentclass[../../../include/open-logic-section]{subfiles}

\begin{document}

\olfileid{sth}{ord-arithmetic}{add}
\olsection{序数加法}

假设我们要将 $\alpha$ 与 $\beta$ 相加。我们只需将 $\beta$ 的\emph{副本}直接置于 $\alpha$ 的副本之后。
（我们需要取\emph{副本}，因为由
\olref[ordinals][basic]{ordinalsaresubsets} 可知，要么 $\alpha
\subseteq \beta$，要么 $\beta \subseteq \alpha$。）
其直观效果是先遍历一个 $\alpha$-序列，\emph{然后再}遍历一个 $\beta$-序列。所得的序列将是良序的；
因此根据
\olref[ordinals][ordtype]{thmOrdinalRepresentation}，它同构于某个（唯一的）序数。该序数即可视为 $\alpha$ 与 $\beta$ 的\emph{和}。

这就是序数加法背后的直观想法。要严格地定义它，我们从取集合\emph{副本}的想法入手。这里的做法是使用任意的标签 $0$ 和 $1$，来标明对象的具体来源：

\begin{defn}\ollabel{defdissum}
$A$ 与 $B$ 的\emph{不交并}是 $A \disjointsum B = (A\times
\{0\}) \cup (B \times \{1\})$。
\end{defn}

我们接下来在序数对上定义一个序：

\begin{defn}
对任意序数 $\alpha_1, \alpha_2, \beta_1, \beta_2$，规定：
\begin{align*}
	\tuple{\alpha_1, \alpha_2} \rlexless \tuple{\beta_1, \beta_2}\text{ 当且仅当 }&
	\text{要么 $\alpha_2 \in \beta_2$}\\
	& \text{要么两者 $\alpha_2 = \beta_2$ 且 $\alpha_1 \in \beta_1$}
\end{align*}
\end{defn}

这是一个\emph{逆字典序}，因为它是先按第二个元素排序，再按第一个元素排序。回想一下，我们希望将 $\alpha \ordplus \beta$ 定义为一个 $\alpha$ 的副本后跟一个 $\beta$ 的副本所构成的序型。为此，我们规定：

\begin{defn}\ollabel{defordplus}
对任意序数 $\alpha$, $\beta$，它们的和是 $\alpha \ordplus
\beta = \ordtype{\alpha \disjointsum \beta, \rlexless}$。
\end{defn}
\noindent
注意，此处我们略微滥用了记号；严格来讲，应当用“$\Setabs{\tuple{x,y}\in \alpha\disjointsum\beta}{x \rlexless y}$”代替“$\rlexless$”。但为简洁起见，下文我们将继续沿用此滥用记号的做法。

以下结果与
\olref[ordinals][ordtype]{thmOrdinalRepresentation} 一起，确认了我们的定义是良定的：

\begin{lem}\ollabel{ordsumlessiswo}
对任意序数 $\alpha$ 和 $\beta$，$\tuple{\alpha \disjointsum \beta, \rlexless}$ 是一个良序。
\end{lem}

\begin{proof}
显然 $\rlexless$ 在 $\alpha \disjointsum \beta$ 上是连通的。为证明它是良基的，取定非空子集 $X \subseteq \alpha
\disjointsum \beta$。令 $Y$ 为 $X$ 中第二坐标最小的那些元素组成的子集，即 $Y = \Setabs{\tuple{\gamma, i} \in X}{(\forall \tuple{\delta, j} \in X)i \leq j}$。然后选取 $Y$ 中第一坐标最小的元素。
\end{proof}
\noindent
这样，我们便得到了一个优美且明确的序数加法定义。以下是一个并不令人意外的事实（回顾 \olref[z][infinity-again]{defnomega}，$1 = \{0\}$）：
\begin{prop}
对任意序数 $\alpha$，有 $\alpha \ordplus  1 = \ordsucc{\alpha}$。
\end{prop}

\begin{proof}
考虑从 $\ordsucc{\alpha} = \alpha \cup
\{\alpha\}$ 到 $\alpha\disjointsum1 = (\alpha \times \{0\})
\disjointsum (\{0\} \times \{1\})$ 的同构映射 $f$，其由 $f(\gamma) =
\tuple{\gamma, 0}$（其中 $\gamma \in \alpha$）与 $f(\alpha) = \tuple{0,
1}$ 给出。
\end{proof}
\noindent
此外，不难证明加法满足以下递归条件：

\begin{lem}\ollabel{ordadditionrecursion}
对任意序数 $\alpha, \beta$，我们有：
\begin{align*}
	\alpha\ordplus 0 &= \alpha\\
	\alpha \ordplus  (\beta\ordplus 1) &= (\alpha \ordplus  \beta) \ordplus  1\\
	\alpha  \ordplus  \beta &= \supstrict_{\delta < \beta}(\alpha \ordplus  \delta) && \text{若 $\beta $ 是极限序数}
\end{align*}
\end{lem}

\begin{proof}
我们逐项验证；首先：
\begin{align*}
	\alpha \ordplus  0
	& = \ordtype{(\alpha \times \{0\}) \cup (0 \times \{1\}), \rlexless} \\
	&= \ordtype{(\alpha \times \{0\}) \cup \{0\}, \rlexless}\\
	&= \alpha\\
	\alpha \ordplus (\beta \ordplus  1)
	%&= \ordtype{(\alpha\times \{0\}) \cup (\ordtype{\beta \ordplus  1}\times \{1\}), \rlexless} \\
	&= \ordtype{(\alpha\times \{0\}) \cup (\ordsucc{\beta}\times \{1\}), \rlexless} \\
	&= \ordtype{(\alpha\times \{0\}) \cup (\beta \times \{1\}), \rlexless} \ordplus 1\\
	&= (\alpha \ordplus  \beta) \ordplus  1
\end{align*}
现设 $\beta \neq \emptyset$ 为极限序数。若 $\delta < \beta$，则亦有 $\delta\ordplus 1 < \beta$，从而 $\alpha \ordplus  \delta$ 是 $\alpha \ordplus  \beta$ 的真前段。因此 $\alpha
\ordplus  \beta$ 是 $X = \Setabs{\alpha
\ordplus  \delta}{\delta < \beta}$ 的严格上界。进一步，若 $\alpha \leq \gamma <
\alpha \ordplus  \beta$，则显然存在某个 $\delta < \beta$ 使得 $\gamma = \alpha \ordplus
\delta$。所以 $\alpha \ordplus  \beta =
\supstrict_{\delta< \beta}(\alpha\ordplus \delta)$。
\end{proof}

不过，这里有一个引人注目的事实。为了定义序数加法，我们\emph{也可以}简单地使用超限递归定理，直接列出递归方程，正如 \olref{ordadditionrecursion} 所给出的那样（只是用“$\ordsucc{\beta}$”而非“$\beta \ordplus 1$”）。

因此，有两种不同的方式可以用来定义序数上的运算。我们可以\emph{综合地}定义它们，即显式地构造一个良序集并考察其序型；或者我们可以\emph{递归地}定义它们，只需列出递归方程即可。然而，只要处理得当，两者的结果完全相同。因为 \olref[ordinals][ordtype]{thmOrdinalRepresentation} 保证了这些递归方程确定的是\emph{唯一}的序数。

在许多方面，序数算术与自然数加法的表现相似。例如，我们可以证明以下结论：

\begin{lem}\ollabel{ordinaladditionisnice}
若 $\alpha, \beta, \gamma$ 是序数，则：
\begin{enumerate}
	\item\ollabel{ordaddition1} 若 $\beta < \gamma$，则 $\alpha
	\ordplus  \beta < \alpha \ordplus  \gamma$
	\item\ollabel{ordaddition2} 若 $\alpha \ordplus  \beta =
	\alpha\ordplus \gamma$，则 $\beta = \gamma$
	\item\ollabel{ordaddition3}  $\alpha \ordplus  (\beta \ordplus
	\gamma) = (\alpha \ordplus  \beta) \ordplus  \gamma$，即加法满足结合律
	\item\ollabel{ordaddition4}  若 $\alpha \leq \beta$，则 $\alpha
	\ordplus  \gamma \leq \beta \ordplus \gamma$
\end{enumerate}
\end{lem}

\begin{proof}
我们证明 \olref{ordaddition3}，其余部分留作练习。
证明对 $\gamma$ 进行简单超限归纳，并使用
\olref{ordadditionrecursion}。
当 $\gamma = 0$ 时：
\[
(\alpha \ordplus  \beta) \ordplus  0 = \alpha \ordplus  \beta  = \alpha \ordplus  (\beta \ordplus  0)
\]
当 $\gamma = \delta\ordplus 1$ 时，归纳假设 $(\alpha
\ordplus  \beta) \ordplus  \delta = \alpha \ordplus  (\beta \ordplus
\delta)$；现三次应用 \olref{ordadditionrecursion}：
\begin{align*}
	(\alpha \ordplus  \beta) \ordplus  (\delta \ordplus  1) & = ((\alpha \ordplus  \beta) \ordplus  \delta)\ordplus 1\\
	& = (\alpha \ordplus  (\beta \ordplus  \delta)) \ordplus  1\\
	& = \alpha \ordplus  ((\beta \ordplus  \delta)\ordplus 1)\\
	& = \alpha \ordplus  (\beta \ordplus  (\delta\ordplus 1))
\end{align*}
当 $\gamma$ 是极限序数时，归纳假设若
$\delta \in \gamma$ 则 $(\alpha \ordplus  \beta) \ordplus  \delta =
\alpha \ordplus  (\beta \ordplus  \delta)$；现：
\begin{align*}
	(\alpha \ordplus  \beta) \ordplus  \gamma & = \supstrict_{\delta < \gamma}((\alpha \ordplus  \beta) \ordplus  \delta) \\
	&= \supstrict_{\delta < \gamma}(\alpha \ordplus  (\beta \ordplus  \delta))\\
	&= \alpha \ordplus  \supstrict_{\delta < \gamma}(\beta \ordplus  \delta)\\
	& = \alpha \ordplus  (\beta \ordplus  \gamma)
\end{align*}
\end{proof}

\begin{prob}
证明
\olref[sth][ord-arithmetic][add]{ordinaladditionisnice} 的其余部分。
\end{prob}

就这些方面而言，序数加法应当让人感到十分熟悉。但在一个关键方面，序数加法\emph{并不}类似于自然数上的加法。

\begin{prop}\ollabel{ordsumnotcommute}
序数加法\emph{不}满足交换律；$1 \ordplus  \omega = \omega <
\omega \ordplus  1$。
\end{prop}

\begin{proof}
注意到 $1 \ordplus  \omega = \supstrict_{n < \omega} (1 \ordplus n)
= \omega \in \omega \cup \{\omega\} = \ordsucc{\omega} = \omega
\ordplus  1$。
\end{proof}
\noindent
虽然这起初可能令人惊讶，但其实\emph{不然}。一方面，当你考虑 $1 \ordplus  \omega$ 时，你思考的是在所有自然数\emph{之前}放置一个额外元素所得到的序型。正如我们在
\olref[sfr][infinite][hilbert]{sec} 中对Hilbert旅馆所作的推理那样，直观上看，这个额外的首元素不应改变整体的序型。另一方面，当你考虑 $\omega \ordplus  1$ 时，你思考的是在所有自然数\emph{之后}放置一个额外元素所得到的序型。那就截然不同了！

\end{document}